% Preambel mit Einstellungen importieren
% Document type and used packages
\documentclass[openany,
%    open=right, % Kapitel darf nur auf rechten Seite beginnen
    paper=A4,               % DIN-A4-Papier
    a4paper,                % DIN-A4-Papier
    12pt,                   % Schriftgöße
    headings=small,         % Kleine Überschriften
    headsepline=true,       % Trennlinie am Kopf der Seite
    footsepline=false,      % Keine Trennlinie am Fuß der Seite
    bibliography=totoc,     % Literaturverzeichnis in das Inhaltsverzeichnis aufnehmen
    twoside=on,             % Doppelseitiger Druck - auf off stellen für einseitig
    DIV=7,                  % Verhältnis der Ränder zum bedruckten Bereich
    chapterprefix=true,     % Kapitel x vor dem Kapitelnamen
    cleardoublepage=plain]{scrbook}

% Pakete einbinden, die benötigt werden
\usepackage{scrpage2}
\usepackage[utf8]{inputenc}       % Dateien in UTF-8 benutzen
\usepackage[T1]{fontenc}          % Zeichenkodierung
\usepackage{graphicx}             % Bilder einbinden
\usepackage[main=ngerman, english]{babel}       % Deutsch und Englisch unterstützen
\usepackage{xcolor}               % Color support
\usepackage{amsmath}              % Matheamtische Formeln
\usepackage{amsfonts}             % Mathematische Zeichensätze
\usepackage{amssymb}              % Mathematische Symbole
\usepackage{float}                % Fließende Objekte (Tabellen, Grafiken etc.)
\usepackage{booktabs}             % Korrekter Tabellensatz
\usepackage[printonlyused]{acronym}  % Abkürzungsverzeichnis [nur verwendete Abkürzugen]
\usepackage{makeidx}              % Sachregister
\usepackage{listings}             % Source Code listings
\usepackage{listingsutf8}         % Listings in UTF8
\usepackage[hang,font={sf,footnotesize},labelfont={footnotesize,bf}]{caption} % Beschriftungen
\usepackage[scaled]{helvet}       % Schrift Helvetia laden
\usepackage[absolute]{textpos}	  % Absolute Textpositionen (für Deckblatt)
\usepackage{calc}                 % Berechnung von Positionen
\usepackage{blindtext}            % Blindtexte
\usepackage[bottom=40mm,left=35mm,right=35mm,top=30mm]{geometry} % Ränder ändern
\usepackage{setspace}             % Abstände korrigieren
\usepackage{ifthen}               % Logische Bedingungen mit ifthenelse
\usepackage{scrhack}              % Get rid of tocbasic warnings
\usepackage[pagebackref=false,german]{hyperref}  % Hyperlinks
\usepackage[all]{hypcap}          % Korrekte Verlinkung von Floats
\usepackage[autostyle=true,german=quotes]{csquotes}   % Zitate
\usepackage[backend=biber,
  isbn=false,                     % ISBN nicht anzeigen, gleiches geht mit nahezu allen anderen Feldern
  sortlocale=de_DE,               % Sortierung der Einträge für Deutsch
  %sortlocale=en_US,              % Sortierung der Einträge für Englisch
  autocite=inline,                % regelt Aussehen für \autocite (inline=\parancite)
  hyperref=true,                  % Hyperlinks für Ziate
  %style=ieee                     % Zitate als Zahlen [1]
  %style=alphabetic               % Zitate als Kürzel und Jahr [Ein05]
  style=authoryear                % Zitate Author und Jahr [Einstein (1905)]
]{biblatex}                       % Literaturverwaltung mit BibLaTeX
\usepackage{rotating}             % Seiten drehen
\usepackage{harveyballs}          % Harveyballs

\usepackage{tikz}                 % Graph

\usepackage{verbatim}
\usetikzlibrary{arrows,shapes}

\usepackage{subcaption}
\DeclareCaptionFont{tiny}{\tiny}
\captionsetup{font+=footnotesize}
\captionsetup[sub]{font+=footnotesize}
\captionsetup[table]{belowskip={10pt}}

\setlength{\bibitemsep}{1em}     % Abstand zwischen den Literaturangaben
\setlength{\bibhang}{2em}        % Einzug nach jeweils erster Zeile

% Trennung von URLs im Literaturverzeichnis (große Werte [> 10000] verhindern die Trennung)
\defcounter{biburlnumpenalty}{10} % Strafe für Trennung in URL nach Zahl
\defcounter{biburlucpenalty}{500}  % Strafe für Trennung in URL nach Großbuchstaben
\defcounter{biburllcpenalty}{500}  % Strafe für Trennung in URL nach Kleinbuchstaben

% Farben definieren
\definecolor{linkblue}{RGB}{0, 0, 100}
\definecolor{linkblack}{RGB}{0, 0, 0}
\definecolor{comment}{RGB}{63, 127, 95}
\definecolor{darkgreen}{RGB}{14, 144, 102}
\definecolor{darkblue}{RGB}{0,0,168}
\definecolor{darkred}{RGB}{128,0,0}
\definecolor{javadoccomment}{RGB}{0,0,240}

% Einstellungen für das Hyperlink-Paket
\hypersetup{
    colorlinks=true,      % Farbige links verwenden
%    allcolors=linkblue,
    linktoc=all,          % Links im Inhaltsverzeichnis
    linkcolor=linkblack,  % Querverweise
    citecolor=linkblack,  % Literaturangaben
	filecolor=linkblack,  % Dateilinks
	urlcolor=linkblack    % URLs
}

% Einstellungen für Quelltexte
\lstset{
      xleftmargin=0.2cm,
      basicstyle=\footnotesize\ttfamily,
      keywordstyle=\color{darkgreen},
      identifierstyle=\color{darkblue},
      commentstyle=\color{comment},
      stringstyle=\color{darkred},
      tabsize=2,
      lineskip={2pt},
      columns=flexible,
      inputencoding=utf8,
      captionpos=b,
      breakautoindent=true,
	  breakindent=2em,
	  breaklines=true,
	  prebreak=,
	  postbreak=,
      numbers=none,
      numberstyle=\tiny,
      showspaces=false,      % Keine Leerzeichensymbole
      showtabs=false,        % Keine Tabsymbole
      showstringspaces=false,% Leerzeichen in Strings
      morecomment=[s][\color{javadoccomment}]{/**}{*/},
      literate={Ö}{{\"O}}1 {Ä}{{\"A}}1 {Ü}{{\"U}}1 {ß}{{\ss}}2 {ü}{{\"u}}1 {ä}{{\"a}}1 {ö}{{\"o}}1
}

\urlstyle{same}

% Einstellungen für Überschriften
\renewcommand*{\chapterformat}{%
  \Large\chapapp~\thechapter   % Große Schrift
  \vspace{0.3cm}               % Abstand zum Titel des Kapitels
}

% Abstände für die Überschriften setzen
\renewcommand{\chapterheadstartvskip}{\vspace*{2.6cm}}
\renewcommand{\chapterheadendvskip}{\vspace*{1.5cm}}

\RedeclareSectionCommand[
  beforeskip=-1.8\baselineskip,
  afterskip=0.25\baselineskip]{section}

\RedeclareSectionCommand[
  beforeskip=-1.8\baselineskip,
  afterskip=0.15\baselineskip]{subsection}

\RedeclareSectionCommand[
  beforeskip=-1.8\baselineskip,
  afterskip=0.15\baselineskip]{subsubsection}


% In der Kopfzeile nur die kurze Kapitelbezeichnung (ohne Kapitel davor)
\renewcommand*\chaptermarkformat{\thechapter\autodot\enskip}
\automark[chapter]{chapter}

% Einstellungen für Schriftarten
\setkomafont{pagehead}{\normalfont\sffamily}
\setkomafont{pagenumber}{\normalfont\sffamily}
\setkomafont{paragraph}{\sffamily\bfseries\small}
\setkomafont{subsubsection}{\sffamily\itshape\bfseries\small}
\addtokomafont{footnote}{\footnotesize}
\setkomafont{chapter}{\LARGE\selectfont\bfseries}

% Wichtige Abstände
\setlength{\parskip}{0.2cm}  % 2mm Abstand zwischen zwei Absätzen
\setlength{\parindent}{0mm}  % Absätze nicht einziehen
\clubpenalty = 10000         % Keine "Schusterjungen"
\widowpenalty = 10000        % Keine "Hurenkinder"
\displaywidowpenalty = 10000 % Keine "Hurenkinder"
\renewcommand{\footnotesize}{\fontsize{9}{10}\selectfont} % Größe der Fußnoten
\setlength{\footnotesep}{8pt} % Abstand zwischen den Fußnoten

% Index erzeugen
\makeindex

% Einfacher Font-Wechsel über dieses Makro
\newcommand{\changefont}[3]{
\fontfamily{#1} \fontseries{#2} \fontshape{#3} \selectfont}

% Eigenes Makro für Bilder
\newcommand{\bild}[3]{
\begin{figure}[h]
  \centering
  \includegraphics[width=#2]{#1}
  \caption{#3}
  \label{#1}
\end{figure}}

% Wo liegt Sourcecode?
\newcommand{\srcloc}{src/}

% Wo sind die Bilder?
\graphicspath{{bilder/}}

% Makros für typographisch korrekte Abkürzungen
\newcommand{\zb}[0]{z.\,B.\ }
\newcommand{\dahe}[0]{d.\,h.\ }
\newcommand{\ua}[0]{u.\,a.\ }

% Flags für Veröffentlichung und Sperrvermerk
\newboolean{hsmapublizieren}
\newboolean{hsmasperrvermerk}


% Dokumenteninfos importieren
% -------------------------------------------------------
% Daten für die Arbeit
% Wenn hier alles korrekt eingetragen wurde, wird das Titelblatt
% automatisch generiert. D.h. die Datei titelblatt.tex muss nicht mehr
% angepasst werden.

\newcommand{\hsmasprache}{de} % de oder en für Deutsch oder Englisch
% Für korrekt sortierte Literatureinträge, noch preambel.tex anpassen
% und zwar bei \usepackage[main=ngerman, english]{babel},
% \usepackage[pagebackref=false,german]{hyperref}
% und \usepackage[autostyle=true,german=quotes]{csquotes}

% Titel der Arbeit auf Deutsch
\newcommand{\hsmatitelde}{Kartierung eines Stockwerkes mittels G-Mapping}

% Titel der Arbeit auf Englisch
\newcommand{\hsmatitelen}{Application of a flux compensator for timetravel with a maximum velocity of warp~7}

% Weitere Informationen zur Arbeit
\newcommand{\hsmaort}{Mannheim}    % Ort
\newcommand{\hsmaautorvname}{Max} % Vorname(n)
\newcommand{\hsmaautornname}{Mustermann} % Nachname(n)
\newcommand{\hsmadatum}{17.07.18} % Datum der Abgabe
\newcommand{\hsmajahr}{2018} % Jahr der Abgabe
\newcommand{\hsmafirma}{Paukenschlag GmbH, Mannheim} % Firma bei der die Arbeit durchgeführt wurde
\newcommand{\hsmabetreuer}{Prof. Jörn Fischer, Hochschule Mannheim} % Betreuer an der Hochschule
\newcommand{\hsmazweitkorrektor}{Prof. Thomas Ihme, Hochschule Mannheim} % Betreuer im Unternehmen oder Zweitkorrektor
\newcommand{\hsmafakultaet}{I} % I für Informatik
\newcommand{\hsmastudiengang}{IB} % IB IMB UIB IM MTB

% Zustimmung zur Veröffentlichung
\setboolean{hsmapublizieren}{false}   % Einer Veröffentlichung wird zugestimmt
\setboolean{hsmasperrvermerk}{false} % Die Arbeit hat keinen Sperrvermerk

% -------------------------------------------------------
% Abstract

% Kurze (maximal halbseitige) Beschreibung, worum es in der Arbeit geht auf Deutsch
\newcommand{\hsmaabstractde}{Abstract}

% Kurze (maximal halbseitige) Beschreibung, worum es in der Arbeit geht auf Englisch

\newcommand{\hsmaabstracten}{The European languages are members of the same family. Their separate existence is a myth. For science, music, sport, etc, Europe uses the same vocabulary. The languages only differ in their grammar, their pronunciation and their most common words. Everyone realizes why a new common language would be desirable: one could refuse to pay expensive translators. To achieve this, it would be necessary to have uniform grammar, pronunciation and more common words. If several languages coalesce, the grammar of the resulting language is more simple and regular than that of the individual languages. The new common language will be more simple and regular than the existing European languages. It will be as simple as Occidental; in fact, it will be Occidental. To an English person, it will seem like simplified English, as a skeptical Cambridge friend of mine told me what Occidental is.}


% Literatur-Datenbank
\addbibresource{literatur.bib}   % BibLaTeX-Datei mit Literaturquellen einbinden

\begin{document}
\frontmatter

% Römische Ziffern für die "Front-Matter"
\setcounter{page}{0}
\changefont{ptm}{m}{n}  % Times New Roman für den Fließtext
\renewcommand{\rmdefault}{ptm}

% Titelblatt
% -------------------------------------------------------
% In dieser Datei sollten eigentlich keine Veränderungen mehr
% notwendig sein.
% -------------------------------------------------------

\thispagestyle{empty}

% Fakultäten der HS-Mannheim
% -------------------------------------------------------
\ifthenelse{\equal{\hsmafakultaet}{I}}%
  {\newcommand{\hsmafakultaetlangde}{Fakultät für Informatik}%
   \newcommand{\hsmafakultaetlangen}{Department of Computer Science}}{}

\ifthenelse{\equal{\hsmafakultaet}{E}}%
  {\newcommand{\hsmafakultaetlangde}{Fakultät für Elektrotechnik}%
   \newcommand{\hsmafakultaetlangen}{Department of Electrical Engineering}}{}

\ifthenelse{\equal{\hsmafakultaet}{S}}%
  {\newcommand{\hsmafakultaetlangde}{Fakultät für Sozialwesen}%
   \newcommand{\hsmafakultaetlangen}{Department of Social Work}}{}
   
\ifthenelse{\equal{\hsmafakultaet}{B}}%
  {\newcommand{\hsmafakultaetlangde}{Fakultät für Biotechnologie}%
   \newcommand{\hsmafakultaetlangen}{Department of Biotechnology}}{}

\ifthenelse{\equal{\hsmafakultaet}{D}}%
  {\newcommand{\hsmafakultaetlangde}{Fakultät für Gestaltung}%
   \newcommand{\hsmafakultaetlangen}{Department of Design}}{}

\ifthenelse{\equal{\hsmafakultaet}{M}}%
  {\newcommand{\hsmafakultaetlangde}{Fakultät für Maschinenbau}%
   \newcommand{\hsmafakultaetlangen}{Department of Mechanical Engineering}}{}

\ifthenelse{\equal{\hsmafakultaet}{N}}%
  {\newcommand{\hsmafakultaetlangde}{Fakultät für Informationstechnik}%
   \newcommand{\hsmafakultaetlangen}{Department of Information Technology}}{}
   
\ifthenelse{\equal{\hsmafakultaet}{W}}%
  {\newcommand{\hsmafakultaetlangde}{Fakultät für Wirtschaftsingenieurwesen}%
   \newcommand{\hsmafakultaetlangen}{Department of Engineering and Management}}{}
   
\ifthenelse{\equal{\hsmafakultaet}{C}}%
  {\newcommand{\hsmafakultaetlangde}{Fakultät für Verfahrens- und Chemietechnik}%
   \newcommand{\hsmafakultaetlangen}{Department of Chemical Process Engineering}}{}
   

\ifthenelse{\equal{\hsmastudiengang}{IB}}%
  {\newcommand{\hsmastudienganglangde}{Informatik}%
  \newcommand{\hsmastudienganglangen}{Computer Science}%
  \newcommand{\hsmatypde}{Ausarbeitung}%
  \newcommand{\hsmatypen}{Bachelor Thesis}%
  \newcommand{\hsmagrad}{\hsmabsc}}{}

\ifthenelse{\equal{\hsmastudiengang}{IMB}}%
  {\newcommand{\hsmastudienganglangde}{Medizinische Informatik}%
  \newcommand{\hsmastudienganglangen}{Medical Informatics}%
  \newcommand{\hsmatypde}{Bachelor-Thesis}%
  \newcommand{\hsmatypen}{Bachelor Thesis}%
  \newcommand{\hsmagrad}{\hsmabsc}}{}
  
\ifthenelse{\equal{\hsmastudiengang}{UIB}}%
  {\newcommand{\hsmastudienganglangde}{Unternehmens- und Wirtschaftsinformatik}%
  \newcommand{\hsmastudienganglangen}{Enterprise Computing}%  
  \newcommand{\hsmatypde}{Bachelor-Thesis}%
  \newcommand{\hsmatypen}{Bachelor Thesis}%
  \newcommand{\hsmagrad}{\hsmabsc}}{}

\ifthenelse{\equal{\hsmastudiengang}{IM}}%
  {\newcommand{\hsmastudienganglangde}{Informatik}%
   \newcommand{\hsmastudienganglangen}{Computer Science}%
   \newcommand{\hsmatypde}{Master-Thesis}%
   \newcommand{\hsmatypen}{Master Thesis}%
   \newcommand{\hsmagrad}{\hsmamaster}}{}

\ifthenelse{\equal{\hsmastudiengang}{MEB}}%
  {\newcommand{\hsmastudienganglangde}{Mechatronik}%
   \newcommand{\hsmastudienganglangen}{Mechatronic}%
   \newcommand{\hsmatypde}{Bachelor-Thesis}%
   \newcommand{\hsmatypen}{Bachelor Thesis}%
   \newcommand{\hsmagrad}{\hsmabsc}}{}
   
\ifthenelse{\equal{\hsmastudiengang}{UB}}%
  {\newcommand{\hsmastudienganglangde}{Automatisierungstechnik}%
   \newcommand{\hsmastudienganglangen}{Automation Technology}%
   \newcommand{\hsmatypde}{Bachelor-Thesis}%
   \newcommand{\hsmatypen}{Bachelor Thesis}%
   \newcommand{\hsmagrad}{\hsmabsc}}{}
   
\ifthenelse{\equal{\hsmastudiengang}{ELB}}%
  {\newcommand{\hsmastudienganglangde}{Elektro- und Informationstechnik/Ingenieurpädagogik}%
   \newcommand{\hsmastudienganglangen}{Elektro- und Informationstechnik/Ingenieurpädagogik}%
   \newcommand{\hsmatypde}{Bachelor-Thesis}%
   \newcommand{\hsmatypen}{Bachelor Thesis}%
   \newcommand{\hsmagrad}{\hsmabsc}}{} 
   
\ifthenelse{\equal{\hsmastudiengang}{EBE}}%
  {\newcommand{\hsmastudienganglangde}{Energietechnik und erneuerbare Energien}%
   \newcommand{\hsmastudienganglangen}{Power Engineering ans Renewable Energies}%
   \newcommand{\hsmatypde}{Bachelor-Thesis}%
   \newcommand{\hsmatypen}{Bachelor Thesis}%
   \newcommand{\hsmagrad}{\hsmabsc}}{}

\ifthenelse{\equal{\hsmastudiengang}{TS}}%
  {\newcommand{\hsmastudienganglangde}{Translation Studies}%
   \newcommand{\hsmastudienganglangen}{Translation Studies}%
   \newcommand{\hsmatypde}{Bachelor-Thesis}%
   \newcommand{\hsmatypen}{Bachelor Thesis}%
   \newcommand{\hsmagrad}{\hsmabsc}}{}
  
\ifthenelse{\equal{\hsmastudiengang}{EM}}%
  {\newcommand{\hsmastudienganglangde}{Automatisierungs- und Energiesysteme}%
   \newcommand{\hsmastudienganglangen}{Automation and Energy Systems}%
   \newcommand{\hsmatypde}{Master-Thesis}%
   \newcommand{\hsmatypen}{Master Thesis}%
   \newcommand{\hsmagrad}{\hsmamaster}}{}
   
\ifthenelse{\equal{\hsmastudiengang}{ELM}}%
  {\newcommand{\hsmastudienganglangde}{Lehramt Ingenieurpädagogik}%
   \newcommand{\hsmastudienganglangen}{Lectureship Educational Engineering}%
   \newcommand{\hsmatypde}{Master-Thesis}%
   \newcommand{\hsmatypen}{Master Thesis}%
   \newcommand{\hsmagrad}{\hsmamaster}}{}
   
\ifthenelse{\equal{\hsmastudiengang}{SAB}}%
  {\newcommand{\hsmastudienganglangde}{Soziale Arbeit}%
   \newcommand{\hsmastudienganglangen}{Social Labour}%
   \newcommand{\hsmatypde}{Bachelor-Thesis}%
   \newcommand{\hsmatypen}{Bachelor Thesis}%
   \newcommand{\hsmagrad}{\hsmaba}}{}
   
\ifthenelse{\equal{\hsmastudiengang}{SAM}}%
  {\newcommand{\hsmastudienganglangde}{Soziale Arbeit}%
   \newcommand{\hsmastudienganglangen}{Social Labour}%
   \newcommand{\hsmatypde}{Master-Thesis}%
   \newcommand{\hsmatypen}{Master Thesis}%
   \newcommand{\hsmagrad}{\hsmamastera}}{}
   
\ifthenelse{\equal{\hsmastudiengang}{BB}}%
  {\newcommand{\hsmastudienganglangde}{Biotechnology}%
   \newcommand{\hsmastudienganglangen}{Biotechnology}%
   \newcommand{\hsmatypde}{Bachelor-Thesis}%
   \newcommand{\hsmatypen}{Bachelor Thesis}%
   \newcommand{\hsmagrad}{\hsmabsc}}{}
   
\ifthenelse{\equal{\hsmastudiengang}{BCB}}%
  {\newcommand{\hsmastudienganglangde}{Biologische Chemie}%
   \newcommand{\hsmastudienganglangen}{Biological Chemics}%
   \newcommand{\hsmatypde}{Bachelor-Thesis}%
   \newcommand{\hsmatypen}{Bachelor Thesis}%
   \newcommand{\hsmagrad}{\hsmabsc}}{}
   
\ifthenelse{\equal{\hsmastudiengang}{BMEBST}}%
  {\newcommand{\hsmastudienganglangde}{Biotechnology - Biomedical Science and Technology}%
   \newcommand{\hsmastudienganglangen}{Biotechnology - Biomedical Science and Technology}%
   \newcommand{\hsmatypde}{Master-Thesis}%
   \newcommand{\hsmatypen}{Master Thesis}%
   \newcommand{\hsmagrad}{\hsmamaster}}{}
   
\ifthenelse{\equal{\hsmastudiengang}{BMEBPD}}%
  {\newcommand{\hsmastudienganglangde}{Biotechnology - Bioprocess Development}%
   \newcommand{\hsmastudienganglangen}{Biotechnology - Bioprocess Development}%
   \newcommand{\hsmatypde}{Master-Thesis}%
   \newcommand{\hsmatypen}{Master Thesis}%
   \newcommand{\hsmagrad}{\hsmamaster}}{}
   
\ifthenelse{\equal{\hsmastudiengang}{BLSM}}%
  {\newcommand{\hsmastudienganglangde}{Life Science Management}%
   \newcommand{\hsmastudienganglangen}{Life Science Management}%
   \newcommand{\hsmatypde}{Master-Thesis}%
   \newcommand{\hsmatypen}{Master Thesis}%
   \newcommand{\hsmagrad}{\hsmamaster}}{}
   
\ifthenelse{\equal{\hsmastudiengang}{KDB}}%
  {\newcommand{\hsmastudienganglangde}{Kommunikationsdesign}%
   \newcommand{\hsmastudienganglangen}{Communication Design}%
   \newcommand{\hsmatypde}{Bachelor-Thesis}%
   \newcommand{\hsmatypen}{Bachelor Thesis}%
   \newcommand{\hsmagrad}{\hsmaba}}{}
   
\ifthenelse{\equal{\hsmastudiengang}{KDM}}%
  {\newcommand{\hsmastudienganglangde}{Kommunikationsdesign}%
   \newcommand{\hsmastudienganglangen}{Communication Design}%
   \newcommand{\hsmatypde}{Master-Thesis}%
   \newcommand{\hsmatypen}{Master Thesis}%
   \newcommand{\hsmagrad}{\hsmamastera}}{}
   
\ifthenelse{\equal{\hsmastudiengang}{MB}}%
  {\newcommand{\hsmastudienganglangde}{Maschinenbau}%
   \newcommand{\hsmastudienganglangen}{Mechanical Engineering}%
   \newcommand{\hsmatypde}{Bachelor-Thesis}%
   \newcommand{\hsmatypen}{Bachelor Thesis}%
   \newcommand{\hsmagrad}{\hsmabsc}}{}
   
\ifthenelse{\equal{\hsmastudiengang}{MM}}%
  {\newcommand{\hsmastudienganglangde}{Maschinenbau}%
   \newcommand{\hsmastudienganglangen}{Mechanical Engineering}%
   \newcommand{\hsmatypde}{Master-Thesis}%
   \newcommand{\hsmatypen}{Master Thesis}%
   \newcommand{\hsmagrad}{\hsmamaster}}{}
   
\ifthenelse{\equal{\hsmastudiengang}{NEB}}%
  {\newcommand{\hsmastudienganglangde}{Elektronik}%
   \newcommand{\hsmastudienganglangen}{Electronics}%
   \newcommand{\hsmatypde}{Bachelor-Thesis}%
   \newcommand{\hsmatypen}{Bachelor Thesis}%
   \newcommand{\hsmagrad}{\hsmabsc}}{}
   
\ifthenelse{\equal{\hsmastudiengang}{TIB}}%
  {\newcommand{\hsmastudienganglangde}{Technische Informatik}%
   \newcommand{\hsmastudienganglangen}{Technical Information Technology}%
   \newcommand{\hsmatypde}{Bachelor-Thesis}%
   \newcommand{\hsmatypen}{Bachelor Thesis}%
   \newcommand{\hsmagrad}{\hsmabsc}}{}
   
\ifthenelse{\equal{\hsmastudiengang}{MTB}}%
  {\newcommand{\hsmastudienganglangde}{Medizintechnik}%
   \newcommand{\hsmastudienganglangen}{Medical Technology}%
   \newcommand{\hsmatypde}{Bachelor-Thesis}%
   \newcommand{\hsmatypen}{Bachelor Thesis}%
   \newcommand{\hsmagrad}{\hsmabsc}}{}
   
\ifthenelse{\equal{\hsmastudiengang}{MTM}}%
  {\newcommand{\hsmastudienganglangde}{Medizintechnik}%
   \newcommand{\hsmastudienganglangen}{Medical Technology}%
   \newcommand{\hsmatypde}{Master-Thesis}%
   \newcommand{\hsmatypen}{Master Thesis}%
   \newcommand{\hsmagrad}{\hsmamaster}}{}
   
\ifthenelse{\equal{\hsmastudiengang}{NM}}%
  {\newcommand{\hsmastudienganglangde}{Informationstechnik}%
   \newcommand{\hsmastudienganglangen}{Informationstechnik}%
   \newcommand{\hsmatypde}{Master-Thesis}%
   \newcommand{\hsmatypen}{Master Thesis}%
   \newcommand{\hsmagrad}{\hsmamaster}}{}
   
\ifthenelse{\equal{\hsmastudiengang}{WB}}%
  {\newcommand{\hsmastudienganglangde}{Wirtschaftsingenieurwesen}%
   \newcommand{\hsmastudienganglangen}{Business Administration and Engineering}%
   \newcommand{\hsmatypde}{Bachelor-Thesis}%
   \newcommand{\hsmatypen}{Bachelor Thesis}%
   \newcommand{\hsmagrad}{\hsmabsc}}{}
   
\ifthenelse{\equal{\hsmastudiengang}{WM}}%
  {\newcommand{\hsmastudienganglangde}{Wirtschaftsingenieurwesen}%
   \newcommand{\hsmastudienganglangen}{Business Administration and Engineering}%
   \newcommand{\hsmatypde}{Master-Thesis}%
   \newcommand{\hsmatypen}{Master Thesis}%
   \newcommand{\hsmagrad}{\hsmamaster}}{}
   
\ifthenelse{\equal{\hsmastudiengang}{VB}}%
  {\newcommand{\hsmastudienganglangde}{Verfahrenstechnik}%
   \newcommand{\hsmastudienganglangen}{Process Engineering}%
   \newcommand{\hsmatypde}{Bachelor-Thesis}%
   \newcommand{\hsmatypen}{Bachelor Thesis}%
   \newcommand{\hsmagrad}{\hsmabsc}}{}
   
\ifthenelse{\equal{\hsmastudiengang}{CB}}%
  {\newcommand{\hsmastudienganglangde}{Chemische Technik}%
   \newcommand{\hsmastudienganglangen}{Chemical Engineering}%
   \newcommand{\hsmatypde}{Bachelor-Thesis}%
   \newcommand{\hsmatypen}{Bachelor Thesis}%
   \newcommand{\hsmagrad}{\hsmabsc}}{}
   
\ifthenelse{\equal{\hsmastudiengang}{CM}}%
  {\newcommand{\hsmastudienganglangde}{Chemieingenieurwesen}%
   \newcommand{\hsmastudienganglangen}{Chemical Engineering}%
   \newcommand{\hsmatypde}{Master-Thesis}%
   \newcommand{\hsmatypen}{Master Thesis}%
   \newcommand{\hsmagrad}{\hsmamaster}}{}

\newcommand{\hsmabsc}{Bachelor of Science (B.Sc.)}
\newcommand{\hsmaba}{Bachelor of Arts (B.A.)}
\newcommand{\hsmamaster}{Master of Science (M.Sc.)}
\newcommand{\hsmamastera}{Master of Arts (M.A.)}
\newcommand{\hsmamasterba}{Master of Business Administration (MBA)}

\newcommand{\hsmakoerperschaftde}{Hochschule Mannheim}
\newcommand{\hsmakoerperschaften}{University of Applied Sciences Mannheim}

\newcommand{\hsmaautorbib}{\hsmaautornname, \hsmaautorvname} % Autor Nachname, Vorname
\newcommand{\hsmaautor}{\hsmaautorvname \ \hsmaautornname} % Autor Vorname Nachname

\ifthenelse{\equal{\hsmasprache}{de}}%
  {\newcommand{\hsmatyp}{\hsmatypde}%
   \newcommand{\hsmathesistype}{in Wahlfachmodul Robotik}%
   \newcommand{\hsmakoerperschaft}{\hsmakoerperschaftde}%
   \newcommand{\hsmastudiengangname}{Studiengang \hsmastudienganglangde}%
   \newcommand{\hsmastudienganglang}{\hsmastudienganglangde}%
   \newcommand{\hsmatitel}{\hsmatitelde}%
   \newcommand{\hsmatutor}{Professor}%
   \newcommand{\hsmafakultaetlang}{\hsmafakultaetlangde}%
   \newcommand{\hsmalistoftables}{Tabellenverzeichnis}%
   \newcommand{\hsmalistoffigures}{Abbildungsverzeichnis}%
   \newcommand{\hsmalistings}{Quellcodeverzeichnis}%
   \newcommand{\hsmaindex}{Index}%
   \newcommand{\hsmaabbreviations}{Abkürzungsverzeichnis}%   
   \selectlanguage{ngerman}}%
  {\newcommand{\hsmatyp}{\hsmatypen}%
   \newcommand{\hsmathesistype}{for the acquisition of the academic degree \hsmagrad}%
   \newcommand{\hsmakoerperschaft}{\hsmakoerperschaften}%
   \newcommand{\hsmastudiengangname}{Course of Studies: \hsmastudienganglang}%
   \newcommand{\hsmastudienganglang}{\hsmastudienganglangen}%
   \newcommand{\hsmatitel}{\hsmatitelen}%
   \newcommand{\hsmatutor}{Tutors}
   \newcommand{\hsmafakultaetlang}{\hsmafakultaetlangen}%
   \newcommand{\hsmalistoftables}{List of Tables}%
   \newcommand{\hsmalistoffigures}{List of Figures}%
   \newcommand{\hsmalistings}{Listings}%
   \newcommand{\hsmaindex}{Index}%
   \newcommand{\hsmaabbreviations}{List of Abbreviations}%
   \selectlanguage{english}}%


% Daten in die Standard-Felder von KOMA-Script eintragen
\titlehead{\hsmatyp\ in\  \hsmastudienganglang}
\subject{}
\title{\hsmatitel}
\author{Minh Le, \ Nicolas Pilch, \ Stefan Witsuba}
\date{\small{\hsmadatum}}

% Daten für das fertige PDF-Dokument
\hypersetup{
  pdftitle={\hsmatitel},  % Titel des Dokuments
  pdfauthor={\hsmaautor},              % Autor
  pdfsubject={\hsmatyp\ in\ \hsmastudienganglang},                % Thema
  pdfkeywords={\hsmatitel}         % Schlüsselworte
}

\newlength{\bindekorrektur}
\newlength{\seitenanfang}
\newlength{\seitenbreite}
  
\setlength{\bindekorrektur}{-46mm}   % Korrektur der horizontalen Position
\setlength{\seitenanfang}{0mm}       % Korrektur der vertikalen Position
\setlength{\seitenbreite}{297mm}

\noindent\includegraphics[width=7cm]{hsma-logo.pdf}\\

% Titel der Arbeit
\begin{textblock*}{128mm}(45mm,\seitenanfang + 62mm) % 4,5cm vom linken Rand und 6,0cm vom oberen Rand
  \centering\Large\sffamily
  \vspace{4mm} % Kleiner zusätzlicher Abstand oben für bessere Optik
  \textbf{\hsmatitel}
\end{textblock*}%

% Name
\begin{textblock*}{128mm}(45mm,\seitenanfang + 103mm)
  \centering\large\sffamily
  Minh Le, \ Nicolas Pilch, \ Stefan Wistuba
\end{textblock*}

% Thesis
\begin{textblock*}{\seitenbreite}(\bindekorrektur,\seitenanfang + 130mm)
  \centering\large\sffamily
  \hsmatyp\\
  \begin{small}\hsmathesistype \end{small}\\
  \vspace{2mm}
  \hsmastudiengangname
\end{textblock*}

% Fakultät
\begin{textblock*}{\seitenbreite}(\bindekorrektur,\seitenanfang + 165mm)
  \centering\large\sffamily
  \hsmafakultaetlang\\
  \vspace{2mm}
  \hsmakoerperschaft
\end{textblock*}

% Datum
\begin{textblock*}{\seitenbreite}(\bindekorrektur,\seitenanfang + 190mm)
  \centering\large 
  \textsf{\hsmadatum}
\end{textblock*}

% Firma
\begin{textblock*}{\seitenbreite}(\bindekorrektur,\seitenanfang + 215mm)
  \centering\large 
  %\textsf{Durchgeführt bei der Firma \hsmafirma}
\end{textblock*}

% Betreuer
\begin{textblock*}{\seitenbreite}(\bindekorrektur,\seitenanfang + 240mm)
  \centering\large\sffamily
  \hsmatutor \\
  \vspace{2mm}
  \hsmabetreuer\\
  \vspace{2mm}
%  \hsmazweitkorrektor
\end{textblock*}

% % Bibliographische Informationen
% \null\newpage
% \thispagestyle{empty}
  
% \newcommand{\hsmabibde}{\begin{small}\textbf{\hsmaautorbib}: \\ \hsmatitelde \ / \hsmaautor. \ -- \\ \hsmatypde, \hsmaort : \hsmakoerperschaftde, \hsmajahr. \pageref{lastpage} Seiten.\end{small}}

% \newcommand{\hsmabiben}{\begin{small}\textbf{\hsmaautorbib}: \\ \hsmatitelen \ / \hsmaautor. \ -- \\ \hsmatypen, \hsmaort : \hsmakoerperschaften, \hsmajahr. \pageref{lastpage} pages. \end{small}}

% \ifthenelse{\equal{\hsmasprache}{de}}%
%   {\hsmabibde \\ \vspace{0.5cm} \\ \hsmabiben}
%   {\hsmabiben \\ \vspace{0.5cm} \\ \hsmabibde}


% % Erklärung
% \clearpage\setcounter{page}{1}
% \thispagestyle{empty}
% \textsf{\large\textbf{Erklärung}}

% Hiermit erkläre ich, dass ich die vorliegende Arbeit selbstständig verfasst und keine anderen als die angegebenen Quellen und Hilfsmittel benutzt habe.

% \ifthenelse{\boolean{hsmapublizieren} \and \not\boolean{hsmasperrvermerk}}%
% {
% \vspace{0.5cm}
% Ich bin damit einverstanden, dass meine Arbeit veröffentlicht wird, d.\,h. dass die Arbeit elektronisch gespeichert, in andere Formate konvertiert, auf den Servern der Hochschule Mannheim öffentlich zugänglich gemacht und über das Internet verbreitet werden darf. 
% }{}%


% \vspace{1cm}
% \hsmaort, \hsmadatum \\

% \vspace{1.2cm}						                                      
% \hsmaautor

% \ifthenelse{\boolean{hsmasperrvermerk}}%
% {%
% \vspace{11cm}
% \color{red}\textsf{\large\textbf{Sperrvermerk}}

% Diese Arbeit basiert auf internen und vertraulichen Daten des Unternehmens \hsmafirma.

% Diese Arbeit darf Dritten, mit Ausnahme der betreuenden Dozenten und befugten Mitglieder des Prüfungsausschusses, ohne ausdrückliche Zustimmung des Unternehmens und des Verfassers nicht zugänglich gemacht werden.

% Eine Vervielfältigung und Veröffentlichung der Arbeit ohne ausdrückliche Genehmigung -- auch in Auszügen -- ist nicht erlaubt.
% \color{black}
% }{}

% \cleardoublepage

% Abstract
\chapter*{Abstract}

\subsubsection*{\hsmatitelde}\hsmaabstractde

% \ifthenelse{\equal{\hsmasprache}{de}}%
%   {\subsubsection*{\hsmatitelde}\hsmaabstractde\subsubsection*{\hsmatitelen}\hsmaabstracten}
%   {\subsubsection*{\hsmatitelen}\hsmaabstracten\subsubsection*{\hsmatitelde}\hsmaabstractde}


% Inhaltsverzeichnis erzeugen
%\cleardoublepage
\pdfbookmark{\contentsname}{Contents}
\tableofcontents

% Korrigiert Nummerierung bei mehrseitigem Inhaltsverzeichnis
%\cleardoublepage
\newcounter{frontmatterpage}
\setcounter{frontmatterpage}{\value{page}}

% Arabische Zahlen für den Hauptteil
\mainmatter

% Den Hauptteil mit vergrößertem Zeilenabstand setzen
\onehalfspacing

% ------------------------------------------------------------------
% Hauptteil der Arbeit
\chapter{Einleitung} % (fold)
\label{cha:einleitung}

Ein großes Problem für Roboter ist es, dass sie sich in ihrer eigenen Umgebung zurechtfinden und anhand von Merkmalen der Umgebung ihre eigene Position erkennen müssen. Ein gutes Beispiel hierfür sind Staubsaugerroboter. Diese fahren meist von einer Basisstation los und erkunden die Umgebung. Hat der Roboter dann seine Aufgabe erledigt, fährt er wieder zu seiner Basisstation zurück. Um Objekte, wie z.\,B. die Basisstation, in einer Welt wiederfinden zu können, muss diese zunächst kartiert werden. Mit einer möglichen Lösung für die Kartierung einer Umgebung beschäftigt sich diese Arbeit. Konkret geht es dabei um den Pioneer 3-DX Roboter, der mit Hilfe eines Laserscanners die Umgebung kartiert. Im Unterschied zum Staubsaugerroboter geschieht dies allerdings nicht autonom, sondern der Roboter wird mittels eines Joysticks durch die Welt gesteuert. Der Roboter verwendet dabei diverse Pakete des \acl{ROS} um zunächst zu Kartieren und eine Karte der Umgebung zu erhalten.\par
In der vorliegenden Projektarbeit werden die Grundlagen von \ac{ROS}, das Vorgehen bei der Kartierung sowie beim Abspielen fertiger Karten beschrieben. Das Ergebnis der vorliegenden Arbeit sind Anleitungen für die genannten Vorgänge, Informationen über alles was man dabei beachten sollte um Probleme zu vermeiden, sowie Beispiele in Form zwei unterschiedlicher Kartierungen.

% chapter einleitung (end) % Externe Datei einbinden
\chapter{Robotik} % (fold)
\label{cha:robotik}

Robotik

\section{Kartierung} % (fold)
\label{sec:kartierung}

% section kartierung (end)

\section{ROS} % (fold)
\label{sec:ros}


% section ros (end)

\section{Pioneer} % (fold)
\label{sec:pioneer}

% section pioneer (end)

% chapter robotik (end) % Externe Datei einbinden
%\chapter{Kartierung} % (fold)
\label{cha:kartierung}

Falls ein Roboter von seiner Umgebung eine Karte erstellen soll, muss er dafür seine Position lokalisieren. Hierfür bracht er jedoch eine Karte seiner Umgebung. Durch diese wechselseitige Abhängigkeit wird das Problem oft auch als \ac{SLAM} bezeichnet. Ein Verfahren zur Lösung des Problems ist der Einsatz von Partikelfiltern.

\section{GMapping} % (fold)
\label{sec:gmapping}

Aus der Odometrie des Roboters und den Beobachtungen kann näherungsweise die Karte und die Bewegung des Roboters berechnet werden. Der Roboter setzt hierzu Rao-Blackwellized Partikelfilter ein. Zuerst wird die Trajektorie des Roboters und daraus anschließend die Karte bestimmt.

Durch einen Rao-Blackwellized Partikelfilter soll eine Anäherung der Umgebungskarte $m$ und des Bewegungsbahn des Roboters $x_{1:t} = x_1,...,x_t$ erechnet werden. Zur Berechnung werden die Odometrie des Roboters $u_{1:t} = u_1,...,u_{t-1}$ und die Umgebungsbeobachtungen benutzt. Allgemein soll die Wahrscheinlichkeit $p(x_{t:1},m | z_{1:t},u_{1:t-1})$ bestimmt werden. Durch Einsatz der Faktorisierung
\begin{equation}
	$p(x_{t:1},m | z_{1:t},u_{1:t-1}) = p(m|x_{1:t},z_{1:t})*p(x_{1:t}|z_{1:t},u_{1:t-1})$
\end{equation}
kann nun zuerst näherungsweise die Trajektorie des Roboters und dann daraus die Karte berechnet.\par
Während $p(m|x_{1:t},z_{1:t})$ analytisch aus gegebenen $x_{1:t}$ und $z_{1:t}$ berechnet werden kann, wird für $p(x_{1:t}|z_{1:t},u_{1:t-1})$ ein Partikelfilter eingesetzt. 


% section gmapping (end)

% chapter kartierung (end) % Externe Datei einbinden
\chapter{Durchführung} % (fold)
\label{cha:durchfuehrung}

Bei der Erstellung einer Karte mittels des Pioneer 3-DX Roboters gibt es einige Dinge zu beachten, damit der Roboter wie gewünscht funktioniert und die Probleme beim Mapping der Karte verringert werden.\par
Zunächst gilt es aber den Roboter mit allen Funktionen richtig zu starten.
Als erster Schritt wird der auf dem Ubuntu-Rechner angemeldete Nutzer mit den nötigen Rechten für die Benutzung der benötigten Pe­ri­phe­rie­ge­räte ausgestattet.

\begin{lstlisting}[language=sh,caption=Rechteverteilung]
    sudo chmod 777 /dev/ttyS0
    sudo chmod 777 /dev/ttyUSB0
    sudo chmod 777 /dev/input/js0
\end{lstlisting}

Anschließend wird  \ac{ROS} mit der angegebenen launch Datei gestartet. Diese muss während des gesamten Vorganges offen bleiben.

\begin{lstlisting}[language=sh,caption=Starten von \ac{ROS}]
    roslaunch p2os_launch pioneerMasterLaunchFile.launch
\end{lstlisting}

Darauf folgend kann in einem neuen Konsolenfenster eine Message für die Motorsteuerung abgesendet werden.

\begin{lstlisting}[language=sh,caption=Starten des Motors]
    rostopic pub /cmd_motor_state p2os_msgs/MotorState 1
\end{lstlisting}

In einem neuen Konsolenfenster kann nun das GMapping gestartet werden. Dass Fenster darf wie alle anderen Konsolenfenster nicht geschlossen werden.

\begin{lstlisting}[language=sh,caption=Starten von GMapping]
    rosrun gmapping slam_gmapping scan:=scan _particles:=100 _linearUpdate:=0.1 _angularUpdate:=0.03
\end{lstlisting}

Nun ist der Roboter einsatzbereit und mittels des Joysticks steuerbar. Außerdem kann ein Livebild der Sensordaten des Laserscanners in Rviz angezeigt werden. 

In einem neuen Konsolenfenster kann nun die Protokollierung der gesendeten Messages gestartet werden, um bei einem Absturz oder ähnliches die Daten neu abspielen und daraus eine Karte generieren zu können.

\begin{lstlisting}[language=sh,caption=Aufnahme aller versendeten Nachrichten]
    rosbag record -O <bag-datei-name> /base_scan /tf /joy /cmd_vel /scan /pose /scan
\end{lstlisting}

Nach Abschluss der Kartierung oder falls der Rechner abgestürzt ist, kann die Karte auf einem leistungsfähigeren Rechner generiert werden.\par

Zuerst müssen alle noch fehlenden Packages für die Kartenerstellung heruntergeladen und installiert werden. Dies geschieht über die Website von Prof. Ihme.\footnote{\url{https://services.informatik.hs-mannheim.de/~ihme/ros/Kartierung}} Es müssen folgende Packages (siehe Tabelle \ref{tab:mapping-packages}) entpacket werden und in den src-Ordner des \ac{ROS} Workspaces verschoben werden.

\begin{table}[H]
  \caption{Benötigte Packages}
  \label{tab:mapping-packages}
  \renewcommand{\arraystretch}{1.2}
  \centering
  \sffamily
  \begin{footnotesize}
    \begin{tabular}{l}
    \toprule
    \textbf{Package}\\
    \midrule
    openslam{\_}gmapping\\
    p2os\\
    pioneer\\
    slam{\_}gmapping\\
    \bottomrule
    \end{tabular}
  \end{footnotesize}
  \rmfamily
\end{table}

Zuerst müssen die Packages installiert werden.\footnote{Das Konsolenfenster muss sich im Verzeichnis \enquoute{catkin_ws} befinden.} 

\begin{lstlisting}[language=sh,caption=Rechte]
    catkin_make
\end{lstlisting}

In einem neuen Konsolenfenster können nun alle noch benötigten Nodes für die Kartierung mittels GMapping gestartet werden. Dieses Fenster darf wie alle anderen und noch zu öffnenden Fenster nicht vor Beendigung des Vorgangs geschlossen werden.

\begin{lstlisting}[language=sh,caption=Rechte]
    roslaunch slam_gmapping_pioneer1.launch
\end{lstlisting}

Zur visuellen Darstellung der Kartierung kann rviz in einem neuen Konsolenfenster geöffnet werden.

\begin{lstlisting}[language=sh,caption=Rechte]
    rosrun rviz rviz -d rviz_kartierung.rviz
\end{lstlisting}

Es öffnet sich ein rviz Fenster, in dem die Karte angezeigt wird. Dies muss während des gesamten Vorganges offen bleiben. Man kann die Kartenerstellung nun mittels geeigneter Tools wie z.\,B. Kazam filmen.

Abschließend müssen noch die aufgenommenen Nachrichten zur Kartenerstellung abgespielt werden.

\begin{lstlisting}[language=sh,caption=Rechte]
    rosbag play --clock <bag-datei-name>
\end{lstlisting}

Wichtig hierbei ist das \enquote{--clock}, wodurch die Zeit in rviz auf den Startzeitpunkt der Karte gesetzt wird und ein Anzeigen ermöglicht wird. Zu beachten ist, dass das Abspielen in Echtzeit erfolgt und genauso lange benötigt wie das Erstellen der Karte. Außerdem kann es vorkommen, dass die Karte bei zwei verschiedenen Wiedergaben, trotz identischer Konfiguration, anders aussieht und das erneute Mapping im schlimmsten Fall nicht funktioniert. Dies zeigt sich vor allem darin, dass sich Teilgebiete überschneiden. Hier gilt aber auch zu beachten, dass sich diese in manchen Fällen wieder richtig ausrichten können und man den Wiedergabevorgang nicht sofort abbrechen muss.\par
Im Rahmen der Lehrveranstaltung wurde der erste Stock des Informatikgebäudes der Hochschule Mannheim an zwei verschiedenen Tagen kartiert.

% chapter durchfuehrung (end) % Externe Datei einbinden
\chapter{Evaluation} % (fold)
\label{cha:evaluation}

Evaluation

% chapter evaluation (end) % Externe Datei einbinden
\chapter{Fazit} % (fold)
\label{cha:fazit}

Fazit

% chapter fazit (end) % Externe Datei einbinden
% ------------------------------------------------------------------

\label{lastpage}

% Neue Seite
\cleardoublepage

% Backmatter mit normalem Zeilenabstand setzen
\singlespacing

% Römische Ziffern für die "Back-Matter", fortlaufend mit "Front-Matter"
\pagenumbering{roman}
\setcounter{page}{\value{frontmatterpage}}

% Abkürzungsverzeichnis
\addchap{\hsmaabbreviations}
% Die längste Abkürzung kann in die eckigen Klammern
% bei \begin{acronym} geschrieben, um einen häßlichen
% Umbruch zu verhindern
\begin{acronym}[IEEE]
\acro{IPC}{Inter-process Communication}
\acro{ROB}{Robotik}
\acro{ROS}{Robot Operating System}
\acro{RPC}{Remote Procedure Call}
\acro{SLAM}{Simultaneous Localization And Mapping}
\end{acronym}


% Tabellenverzeichnis erzeugen
%\cleardoublepage
\phantomsection
\addcontentsline{toc}{chapter}{\hsmalistoftables}
\listoftables

% Abbildungsverzeichnis erzeugen
%\cleardoublepage
\phantomsection
\addcontentsline{toc}{chapter}{\hsmalistoffigures}
\listoffigures

% Listingverzeichnis erzeugen
%\cleardoublepage
% \phantomsection
% \addcontentsline{toc}{chapter}{\hsmalistings}
% \lstlistoflistings

% Literaturverzeichnis erzeugen
\begin{flushleft}
\printbibliography
\end{flushleft}

% Index ausgeben. Wenn Sie keinen Index haben, entfernen Sie einfach
% diesen Teil.
%\cleardoublepage
%\phantomsection
%\addcontentsline{toc}{chapter}{\hsmaindex}
%\printindex

% Anhang. Wenn Sie keinen Anhang haben, entfernen Sie einfach
% diesen Teil.
%\appendix
%\input{kapitel/anhang-a}
%\input{kapitel/anhang-b}

\end{document}
